\documentclass[12pt,reqno]{amsart}
\newcommand{\worksheetnumber}{3.2}
\input{../../101preamble.tex}

\DeclareMathOperator{\Sym}{Sym}
\DeclareMathOperator{\sgn}{sgn}

\title{Worksheet 3.2: Exterior Algebra and Determinants}
\begin{document}
\maketitle
\thispagestyle{firstpage}

Let $d\geq1$. A multilinear map $\beta:V^d\to U$ is \emph{alternating} if $\beta(v_1,\ldots,v_d)=0$ whenever $v_i=v_j$ for some $i\ne j$. For $d=1$ this places no additional condition on a linear map. An alternating map changes sign when two of its inputs are interchanged, as you can see by putting $v+w$ in both of those positions and expanding. The converse needs care. If $2\ne0$ in $\KK$, a map which changes sign under every interchange must vanish whenever two inputs are equal. In characteristic $2$, however, changing the sign does nothing, and a symmetric map need not be alternating. For this reason we will always use the repeated-input condition in our definitions.

You already know an alternating map. On $\RR^2$, the function $((a,b),(c,d))\mapsto ad-bc$ gives the signed area of the parallelogram spanned by two vectors. Reversing their order reverses the sign, and using the same vector twice gives zero. More generally, if $\beta$ is alternating and bilinear, then $\beta(v,w+cv)=\beta(v,w)$, since the extra term is $c\beta(v,v)=0$. Geometrically, sliding one edge of a parallelogram in the direction of the other does not change its signed area. Notice that the formula and this identity make sense over any field, whether or not there is a geometric interpretation.

The goal of this worksheet is to construct a vector space whose linear maps describe alternating maps, just as the tensor product describes multilinear maps, and to use it to define the determinant without choosing a basis. In Worksheet 3.1 we formed the symmetric algebra by imposing commutativity on the tensor algebra. We now impose a different relation: a product with a repeated vector should vanish. The resulting \emph{exterior algebra} describes alternating maps, and its top nonzero degree is a line on which every linear operator acts by a scalar. That scalar is the determinant. We will then use this description to recover several familiar matrix identities, including cofactor expansion, Cramer's rule, and the Cayley--Hamilton theorem.

As in the previous worksheets, all rings have a multiplicative identity but need not be commutative, and the letter $\KK$ will denote a field. All vector spaces, linear maps, and tensor products are over $\KK$ unless we say otherwise, and we will state finite-dimensionality hypotheses where they are needed. We do not assume that the characteristic of $\KK$ is different from $2$. As in Worksheet 3.1, a $\KK$-algebra has an associative, unital, $\KK$-bilinear multiplication.

Recall the tensor algebra $T(V)=\bigoplus_{d\geq0}V^{\otimes d}$ from Worksheet 3.1. Let $J\subset T(V)$ be the two-sided ideal generated by all tensors $v\otimes v$ with $v\in V$. The \emph{exterior algebra} of $V$ is
\[
\Lambda(V)\coloneqq T(V)/J.
\]
The generators of $J$ are homogeneous of degree $2$, so by Worksheet 3.1 the quotient inherits a grading. We write $\Lambda^d(V)$ for its degree $d$ part, called the \emph{$d$th exterior power} of $V$, and
\[
v_1\wedge\cdots\wedge v_d
\]
for the class of $v_1\otimes\cdots\otimes v_d$. Multiplication in $\Lambda(V)$ is called the \emph{wedge product}. In particular, $\Lambda^0(V)=\KK$ and $\Lambda^1(V)=V$, the empty wedge is the multiplicative identity $1\in\KK$, and the wedge of elements of degrees $p$ and $q$ has degree $p+q$.

Notice that the relation $v\wedge v=0$ applies to every vector, and not just to the vectors of a chosen basis. Applying it to $v+w$ gives $v\wedge w=-w\wedge v$. Thus reordering factors changes a sign, and a repeated factor can be moved next to its other occurrence and made to vanish. These observations explain why the quotient describes alternating maps. As with the tensor product, the point is that we have imposed exactly the relations required by the maps we wish to extend. We check this in the following exercise.

\hrulefill

\begin{problems}
\item\label{e:construction} \textbf{Alternating Maps and the Exterior Algebra.} Throughout this exercise let $V$ be a vector space.
\begin{enumerate}
\item Let $\beta:V^d\to U$ be alternating. Fix all inputs except those in positions $i$ and $j$, put $v+w$ in both of these positions, and expand. Prove that interchanging the entries in positions $i$ and $j$ changes the sign of $\beta$. Prove the converse when $2\ne0$ in $\KK$. Over $\FF_2$, show that $(a,b)\mapsto ab$ is unchanged by interchanging its inputs, but is not alternating.
\item Starting from $(v+w)\wedge(v+w)=0$, prove that $v\wedge w=-w\wedge v$. Deduce that the map $V^d\to\Lambda^d(V)$ sending $(v_1,\ldots,v_d)$ to $v_1\wedge\cdots\wedge v_d$ is alternating. Why does this proof also work in characteristic $2$?
\item For $d\geq2$, show that the degree $d$ part $J_d=J\cap V^{\otimes d}$ is spanned by the tensors
\[
v_1\otimes\cdots\otimes v_a\otimes v\otimes v\otimes w_1\otimes\cdots\otimes w_b,\qquad a+b=d-2.
\]
Use the multilinear universal property of $V^{\otimes d}$ from Worksheet 3.1 to prove that every alternating map $\beta:V^d\to U$ induces a unique linear map $\widetilde\beta:\Lambda^d(V)\to U$ with $\widetilde\beta(v_1\wedge\cdots\wedge v_d)=\beta(v_1,\ldots,v_d)$. What gives existence, and what gives uniqueness? What happens when $d=1$?
\item Let $A$ be a $\KK$-algebra, and let $\ell:V\to A$ be a linear map with $\ell(v)^2=0$ for every $v\in V$. Use the universal property of $T(V)$ and the universal property of quotient rings to extend $\ell$ uniquely to an algebra homomorphism $\Lambda(V)\to A$.
\item Let $V=\KK^2$ with basis $e_1,e_2$, and quotient $T(V)$ only by the ideal generated by $e_1\otimes e_1$ and $e_2\otimes e_2$. Show that in degree $2$ the classes of $e_1\otimes e_2$ and $e_2\otimes e_1$ are still linearly independent. Is the square of $e_1+e_2$ zero in this quotient? Explain why the squares of basis vectors alone do not define the exterior algebra.
\end{enumerate}
\end{problems}

\hrulefill

We have constructed $\Lambda(V)$ without choosing a basis. To calculate in it, suppose $V$ is finite dimensional and choose an ordered basis $e_1,\ldots,e_n$. Multilinearity expresses every wedge as a sum of wedges of basis vectors. The relations let us put the indices of each such wedge in increasing order, at the cost of a sign, while a repeated index makes a term vanish. This suggests the following statement.

\begin{theorem}[Basis of an Exterior Power]
If $e_1,\ldots,e_n$ is a basis of $V$, then
\[
\left\{e_{i_1}\wedge\cdots\wedge e_{i_d} \;\; \big| \;\; 1\leq i_1<\cdots<i_d\leq n\right\}
\]
is a basis of $\Lambda^d(V)$. In particular, $\dim(\Lambda^d(V))=\binom{n}{d}$, where $\binom{n}{d}=0$ for $d>n$. For $d=0$, the indicated basis consists of the empty wedge $1$.
\end{theorem}

Spanning follows from the relations, but linear independence needs an argument: passing to a quotient might have imposed further relations among the proposed basis vectors. In Worksheet 2.2 we used coordinate functionals to prove that a tensor basis is linearly independent. Here we modify those functionals by summing over permutations with signs, which produces alternating functions that read off the coefficients of increasing wedges. Notice that this construction does not use determinants, so the basis theorem will not depend on the determinant theory it is about to explain.

It is worth comparing this with symmetric powers. An exterior basis vector uses each $e_i$ either once or not at all, so it is indexed by a subset of the basis, whereas a symmetric monomial allows repetitions. Thus $\Lambda(V)$ has dimension $2^n$ and stops in degree $n$: there is exactly one increasing wedge of all $n$ basis vectors, and none with more than $n$ factors. This one-dimensional top degree will be especially useful, because a linear operator on a line is described by a single scalar. We prove the basis theorem in the following exercise.

\hrulefill

\begin{problems}
\item\label{e:basis} \textbf{A Basis for Exterior Powers.} Throughout this exercise let $e_1,\ldots,e_n$ be a basis of $V$, with dual basis $e_1^{\vee},\ldots,e_n^{\vee}$.
\begin{enumerate}
\item Expand a wedge $v_1\wedge\cdots\wedge v_d$ using the basis of $V$, and then reorder its factors. Prove that the increasing wedges in the theorem span $\Lambda^d(V)$. What happens when $d>n$?
\item For $I=(i_1<\cdots<i_d)$, define
\[
\beta_I(v_1,\ldots,v_d)\coloneqq\sum_{\sigma\in S_d}\sgn(\sigma)\prod_{r=1}^de_{i_r}^{\vee}(v_{\sigma(r)}).
\]
Prove that $\beta_I$ is multilinear. If two inputs are equal, pair each permutation with the permutation obtained by interchanging those two input positions, and prove that the terms cancel. Why is the sum zero even in characteristic $2$? Deduce that $\beta_I$ induces a linear functional on $\Lambda^d(V)$.
\item Evaluate this functional on $e_{j_1}\wedge\cdots\wedge e_{j_d}$ for $j_1<\cdots<j_d$, and show that the value is $1$ if $I=(j_1,\ldots,j_d)$ and $0$ otherwise. Apply these functionals to a linear relation among the increasing wedges, and finish the proof of the theorem.
\item List bases for every nonzero degree of $\Lambda(\KK^3)$. Over $\QQ$, expand
\[
(e_1+2e_2)\wedge(e_2+e_3)\qquad\text{and}\qquad(e_1+e_2)\wedge(e_2+e_3)\wedge(e_3+e_1)
\]
in these bases. What changes when the same expressions are calculated over $\FF_2$?
\item Prove that $v_1\wedge\cdots\wedge v_d\ne0$ if and only if $v_1,\ldots,v_d$ are linearly independent. For one direction, extend the independent vectors to a basis. For the other, express one vector of a dependent list as a linear combination of the others. What happens if the list contains the zero vector?
\end{enumerate}
\end{problems}

\hrulefill

A linear map $\phi:V\to W$ induces a map on exterior powers by applying $\phi$ to every factor:
\[
\Lambda^d(\phi):\Lambda^d(V)\longrightarrow\Lambda^d(W),\qquad v_1\wedge\cdots\wedge v_d\longmapsto\phi(v_1)\wedge\cdots\wedge\phi(v_d).
\]
The universal property of Problem~\ref{e:construction} shows that this is well defined. Taking all degrees together gives a graded algebra homomorphism $\Lambda(\phi):\Lambda(V)\to\Lambda(W)$. Notice that in degree $0$ the map is $\Id_{\KK}$, regardless of $\phi$. As with tensor and symmetric powers, checking composition on wedges of vectors proves functoriality.

The matrix of $\Lambda^d(\phi)$ records how $\phi$ transforms $d$ vectors at once. In degree $1$ we recover the matrix of $\phi$ itself, and in degree $2$ the columns are the expansions of wedges of pairs of columns of that matrix. It is natural to expect that exterior powers detect rank, since there can be no nonzero wedge of more independent vectors in the image than the dimension of the image. We compute some examples, and confirm this expectation, in the following exercise.

\hrulefill

\begin{problems}
\item\label{e:maps} \textbf{Exterior Powers of Linear Maps.} Throughout parts (a) and (d) let $\phi:V\to W$ be a linear map.
\begin{enumerate}
\item Construct $\Lambda^d(\phi)$ using Problem~\ref{e:construction}, and prove that $\Lambda^d(\Id_V)=\Id_{\Lambda^d(V)}$ and $\Lambda^d(\psi\phi)=\Lambda^d(\psi)\Lambda^d(\phi)$. Show that the direct sum of these maps preserves wedge products and the multiplicative identity.
\item Over $\QQ$, let $\phi:\QQ^2\to\QQ^2$ be given by $\phi(e_1)=f_1+f_2$ and $\phi(e_2)=f_2$, where $e_1,e_2$ and $f_1,f_2$ are the standard bases of the source and the target. Compute the matrices of $\Lambda^0(\phi)$, $\Lambda^1(\phi)$, and $\Lambda^2(\phi)$. Compare with your calculation of $\Sym^2(\phi)$ for the same map in Worksheet 3.1.
\item On $\QQ^3$, let $\psi(e_1)=e_1+e_2$, $\psi(e_2)=e_2+e_3$, and $\psi(e_3)=e_1+2e_2+e_3$. Compute the matrix of $\Lambda^2(\psi)$ in the ordered basis $(e_1\wedge e_2,e_1\wedge e_3,e_2\wedge e_3)$, and compute $\Lambda^3(\psi)$. Compare their ranks with $\rank(\psi)$.
\item Suppose $V$ and $W$ are finite dimensional and $\rank(\phi)=r$. Using bases which put $\phi$ in the rank normal form of Worksheet 2.2, prove that
\[
\rank\big(\Lambda^d(\phi)\big)=\binom{r}{d},
\]
including the cases $d=0$ and $d>r$. Explain how to recover $\rank(\phi)$ from knowing which exterior powers of $\phi$ are nonzero.
\end{enumerate}
\end{problems}

\hrulefill

In Worksheet 3.1 the symmetric algebra of a direct sum split as a tensor product of symmetric algebras. There is a corresponding statement for exterior algebras, but now the multiplication on the tensor product must remember a sign. In an exterior algebra, if $a$ and $b$ are homogeneous of degrees $p$ and $q$, then
\[
a\wedge b=(-1)^{pq}\,b\wedge a.
\]
For wedges of vectors, moving one block past the other takes $pq$ interchanges, and the general statement follows by linearity. This property is called \emph{graded commutativity}. Notice that in characteristic $2$ it does not replace the defining relation $v\wedge v=0$.

For graded algebras $A$ and $B$, give $A\otimes B$ the total grading from Worksheet 3.1, with degree $d$ part $\bigoplus_{p+q=d}A_p\otimes B_q$. The \emph{graded tensor product multiplication} is
\[
(a\otimes b)(a'\otimes b')\coloneqq(-1)^{qr}\,(aa')\otimes(bb'),\qquad b\in B_q,\quad a'\in A_r,
\]
on homogeneous elements, extended bilinearly. The sign records moving $b$ past $a'$. We write $A\mathbin{\widehat\otimes}B$ when we want to emphasize this multiplication; the underlying vector space is the ordinary tensor product.

With this convention there is a graded algebra isomorphism, defined without choices,
\[
\Lambda(V\oplus W)\cong\Lambda(V)\mathbin{\widehat\otimes}\Lambda(W).
\]
One direction sends $(v,w)$ to $v\otimes1+1\otimes w$. The other places the factors from $V$ before the factors from $W$ and wedges them together inside $\Lambda(V\oplus W)$. The sign in the multiplication is exactly what makes these two descriptions compatible. For example, in degree $2$ the isomorphism reads
\[
\Lambda^2(V\oplus W)\cong\Lambda^2(V)\oplus(V\otimes W)\oplus\Lambda^2(W).
\]
The three summands account for wedges of two vectors from $V$, of one vector from each space, and of two vectors from $W$. Under our convention a mixed wedge has its $V$ factor first, and writing the $W$ factor first changes its sign. We prove the isomorphism in the following exercise.

\hrulefill

\begin{problems}
\item\label{e:directsum} \textbf{Exterior Powers of a Direct Sum.} Throughout this exercise let $V$ and $W$ be vector spaces.
\begin{enumerate}
\item Prove that $\Lambda(V)$ is graded commutative. For graded algebras $A$ and $B$, check associativity and the identity $1\otimes1$ for the graded tensor product multiplication, by comparing the signs on a product of three homogeneous pure tensors.
\item In $\Lambda(V)\mathbin{\widehat\otimes}\Lambda(W)$, show that $(v\otimes1+1\otimes w)^2=0$ for every $(v,w)\in V\oplus W$. Use Problem~\ref{e:construction}(d) to obtain a graded algebra homomorphism $F:\Lambda(V\oplus W)\to\Lambda(V)\mathbin{\widehat\otimes}\Lambda(W)$.
\item Let $i_V$ and $i_W$ be the coordinate inclusions into $V\oplus W$. Construct a linear map in the other direction by
\[
G(a\otimes b)\coloneqq\Lambda(i_V)(a)\wedge\Lambda(i_W)(b).
\]
Use graded commutativity to prove that $G$ preserves multiplication. Check both composites on the generators $(v,0)$, $(0,w)$, $v\otimes1$, and $1\otimes w$, and deduce that $F$ and $G$ are inverse isomorphisms.
\item Taking degree $d$ parts, obtain an isomorphism
\[
\Lambda^d(V\oplus W)\cong\bigoplus_{p+q=d}\Lambda^p(V)\otimes\Lambda^q(W).
\]
Write down the map from right to left on wedges, and verify that it is compatible with a pair of linear maps $V\to V'$ and $W\to W'$. If $\dim(V)=n$ and $\dim(W)=m$, deduce that $\binom{n+m}{d}=\sum_{p+q=d}\binom{n}{p}\binom{m}{q}$.
\item Take $V=\KK v$ and $W=\KK w$. Compute $(1\otimes w)(v\otimes1)$ and $(v\otimes1)(1\otimes w)$. Why would the ordinary tensor product multiplication give the wrong algebra when the characteristic is different from $2$? What remains true of the squares of the degree $1$ generators in characteristic $2$?
\end{enumerate}
\end{problems}

\hrulefill

Suppose $\dim(V)=n$. By the basis theorem, $\Lambda^n(V)$ is one dimensional, with basis $e_1\wedge\cdots\wedge e_n$ for any basis $e_1,\ldots,e_n$ of $V$. Every endomorphism of a one-dimensional space is multiplication by a unique scalar, and this scalar does not depend on a choice of nonzero vector. We therefore define the \emph{determinant} of a linear map $\phi:V\to V$ by
\[
\Lambda^n(\phi)=\det(\phi)\,\Id_{\Lambda^n(V)}.
\]
Notice that this definition concerns an endomorphism of $V$. For a map between two different $n$-dimensional spaces, the top exterior power is a map between two lines, and turning it into a scalar requires choosing a nonzero vector in each line.

Over $\RR$, this recovers the interpretation of the determinant as a signed volume scaling factor. With the standard basis of $\RR^n$, a wedge $v_1\wedge\cdots\wedge v_n$ is a multiple of $e_1\wedge\cdots\wedge e_n$, and the coefficient is the signed volume of the parallelepiped spanned by $v_1,\ldots,v_n$. The operator multiplies this coefficient by the same factor for every list of vectors. Choosing a different nonzero vector in the top exterior power rescales both the original and the transformed coefficient, leaving their ratio unchanged. The exterior construction retains this comparison over an arbitrary field, where signed Euclidean volume is no longer available.

This definition explains at once why determinants multiply: exterior powers preserve composition, and scalars acting on a line multiply. Choosing a basis recovers the familiar permutation formula. If $\phi(e_j)=\sum_ia_{ij}e_i$, expand $\phi(e_1)\wedge\cdots\wedge\phi(e_n)$. Only the terms which use each basis vector exactly once survive, and reordering such a term introduces the sign of the corresponding permutation. Thus the familiar formula is the coordinate description of an action on a line which was defined without coordinates. We carry out this comparison in the following exercise.

\hrulefill

\begin{problems}
\item\label{e:determinant} \textbf{The Determinant from a Top Exterior Power.} Throughout this exercise let $V$ be finite dimensional, with $\dim(V)=n$.
\begin{enumerate}
\item Show directly that the scalar describing an endomorphism of a one-dimensional space does not change when its basis vector is rescaled. Deduce that $\det(\phi)$ does not depend on a basis of $V$. Explain why the unique endomorphism of the zero space has determinant $1$ under our degree $0$ convention.
\item Use functoriality to prove that $\det(\psi\phi)=\det(\psi)\det(\phi)$ and $\det(\Id_V)=1$. Prove that $\det(\phi)\ne0$ if and only if $\phi$ is invertible, using Problem~\ref{e:basis}(e) or Problem~\ref{e:maps}(d).
\item Let $A=(a_{ij})$ be the matrix of $\phi$. Carry out the wedge expansion described above, and prove that
\[
\det(A)=\sum_{\sigma\in S_n}\sgn(\sigma)\prod_{j=1}^na_{\sigma(j),j}.
\]
Use this formula to prove that $\det(A^{\mathsf T})=\det(A)$, by replacing each permutation with its inverse. Deduce that the determinant is alternating and multilinear in the rows of a matrix, as well as in its columns.
\item Derive the formula for the determinant of a $2\times2$ matrix. Show that adding a multiple of one column to another leaves the determinant unchanged, and that the determinant of a triangular matrix is the product of its diagonal entries. Explain each assertion using either the wedge product or the permutation formula.
\item If $P$ is an invertible matrix, prove that $\det(P^{-1}AP)=\det(A)$. For endomorphisms $\phi$ of $V$ and $\psi$ of $W$, use the isomorphism of Problem~\ref{e:directsum} in top degree to prove that $\det(\phi\oplus\psi)=\det(\phi)\det(\psi)$.
\end{enumerate}
\end{problems}

\hrulefill

The functionals we used to prove the basis theorem can now be described more conceptually. For finite-dimensional $V$, there is a bilinear pairing
\[
\begin{aligned}
\Lambda^d(V^{\vee})\times\Lambda^d(V)&\longrightarrow\KK,\\
(\lambda_1\wedge\cdots\wedge\lambda_d,\ v_1\wedge\cdots\wedge v_d)&\longmapsto\det\big(\lambda_i(v_j)\big)_{1\leq i,j\leq d}.
\end{aligned}
\]
The determinant is alternating in both the functionals and the vectors, so the universal property of exterior powers lets the formula descend in both variables. This pairing is \emph{perfect}, in the sense of Worksheet 3.1: the induced map
\[
\Theta_V:\Lambda^d(V^{\vee})\longrightarrow\big(\Lambda^d(V)\big)^{\vee}
\]
is an isomorphism. Indeed, increasing wedges of a dual basis pair with the corresponding increasing wedges by the identity matrix. Compare this with the symmetric pairing of Worksheet 3.1: this formula needs no division by $d!$, and it is perfect in every characteristic.

The same basis calculation works over a commutative ring. More precisely, let $R$ be a nonzero commutative ring, and let $M$ be a free $R$-module of finite rank $n$, meaning that it has a basis with $n$ elements. Exterior powers over $R$ describe alternating $R$-multilinear maps, and they have the same bases of increasing wedges. With $M^{\vee}\coloneqq\Hom_R(M,R)$, the determinant formula gives an isomorphism
\[
\Lambda_R^d(M^{\vee})\cong\Hom_R\big(\Lambda_R^d(M),R\big).
\]
We will develop tensor products over rings next week. For now, notice why this particular statement does not require division: the pairing of the displayed bases has entries $0$ and $1$. This is a statement about finite free modules. As you saw in Worksheet 2.1, an arbitrary module need not have a basis at all. We check these statements in the following exercise.

\hrulefill

\begin{problems}
\item\label{e:duality} \textbf{Exterior Powers and Duality.} Throughout this exercise let $V$ be finite dimensional.
\begin{enumerate}
\item Starting with the displayed determinant formula, use the universal property of exterior powers twice to construct the pairing. Explain why it is well defined on sums of wedges, as well as on the wedges appearing in the formula.
\item For a basis $e_1,\ldots,e_n$ of $V$ and its dual basis, compute the pairing of two increasing wedges. Prove that $\Theta_V$ is an isomorphism. What happens when $d=0$, where the pairing is multiplication in $\KK$, and when $d>n$?
\item Let $\phi:V\to W$ be a linear map between finite-dimensional spaces. For $\alpha\in\Lambda^d(W^{\vee})$ and $z\in\Lambda^d(V)$, prove that
\[
\big\langle\Lambda^d(\phi^{\vee})(\alpha),z\big\rangle=\big\langle\alpha,\Lambda^d(\phi)(z)\big\rangle.
\]
Translate this equation into the commutativity of the naturality square
\[
\begin{tikzcd}[row sep=4em,column sep=4em]
\Lambda^d(W^{\vee}) \arrow[r,"\Theta_W"] \arrow[d,"\Lambda^d(\phi^{\vee})"'] & (\Lambda^d(W))^{\vee} \arrow[d,"(\Lambda^d(\phi))^{\vee}"]\\
\Lambda^d(V^{\vee}) \arrow[r,"\Theta_V"'] & (\Lambda^d(V))^{\vee}
\end{tikzcd}.
\]
\item On $V=\QQ^3$, take $\alpha=(e_1^{\vee}+e_2^{\vee})\wedge(e_2^{\vee}+e_3^{\vee})$ and $z=(e_1+e_3)\wedge(e_2-e_3)$. Compute $\langle\alpha,z\rangle$ both using the $2\times2$ determinant and by expanding in the exterior bases.
\item Review your calculation in part (b). Which steps use only the commutativity of scalars, and why does the same matrix give an isomorphism for a finite free module over a commutative ring $R$? Compute the pairing between $\Lambda^2_{\ZZ/4\ZZ}\big((\ZZ/4\ZZ)^2\big)$ and the corresponding exterior power of the dual. Does the failure of $2$ to be a unit in $\ZZ/4\ZZ$ cause a problem here?
\end{enumerate}
\end{problems}

\hrulefill

We next compare the exterior description with the cofactor calculations you may have met in linear algebra. Let $A\in\Mat_n(\KK)$ with $n\geq1$, and let $A_{\widehat i,\widehat j}$ be the matrix obtained by deleting row $i$ and column $j$ of $A$. Its determinant is called a \emph{minor} of $A$, and
\[
C_{ij}\coloneqq(-1)^{i+j}\det(A_{\widehat i,\widehat j})
\]
is the corresponding \emph{cofactor}. The \emph{adjugate} of $A$ is the transpose of the matrix of cofactors, $\adj(A)\coloneqq(C_{ij})^{\mathsf T}$, so that $\adj(A)_{ij}=C_{ji}$. When $n=1$, the deleted matrix has size $0\times0$ and determinant $1$.

Cofactors also appear as matrix entries of exterior powers. Relative to a basis $e_1,\ldots,e_n$, put
\[
\omega\coloneqq e_1\wedge\cdots\wedge e_n\qquad\text{and}\qquad\epsilon_i\coloneqq(-1)^{i-1}e_1\wedge\cdots\wedge\widehat{e_i}\wedge\cdots\wedge e_n,
\]
where the hat means that the factor is omitted. Then $e_j\wedge\epsilon_i=\delta_{ij}\omega$. We will see that the matrix of $\Lambda^{n-1}(\phi)$ in the basis $(\epsilon_i)$ is exactly the cofactor matrix of the matrix of $\phi$. Thus the minors of size $n-1$ record the action of $\phi$ on wedges of all but one basis vector. Notice that the adjugate is the \emph{transpose} of this matrix; we will keep track of the transpose carefully in the following exercise.

\begin{theorem}[Adjugate Identity]
For $A\in\Mat_n(\KK)$ with $n\geq1$,
\[
A\adj(A)=\adj(A)A=\det(A)I_n.
\]
The same identities hold for square matrices over any commutative ring.
\end{theorem}

The proof expands determinants along a row or a column. On the diagonal the resulting sum is $\det(A)$, and off the diagonal it is the determinant of a matrix with two equal rows or columns, so it vanishes. No division is involved, which is why the identity holds over a commutative ring as well as over a field. Over a field it gives a formula for the inverse when $\det(A)\ne0$, and this formula leads to Cramer's rule.

For a $2\times2$ matrix, the adjugate is the familiar operation of exchanging the two diagonal entries and negating the other two. In larger dimensions the complementary minors do the same job: multiplying their matrix by $A$ gives a scalar multiple of the identity, and division by that scalar supplies the inverse. Notice that the adjugate is still defined when the determinant is zero. In that case the identity says that every column of $\adj(A)$ lies in $\ker(A)$, so the formula retains information even when it cannot produce an inverse. The exterior interpretation explains why these particular minors occur, while row and column expansions let us verify every entry of the identity. We prove the theorem in the following exercise.

\hrulefill

\begin{problems}
\item\label{e:adjugate} \textbf{Cofactors, the Adjugate, and Cramer's Rule.} Throughout this exercise let $A=(a_{ij})\in\Mat_n(\KK)$ with $n\geq1$, and let $\phi:\KK^n\to\KK^n$ be the linear map with matrix $A$.
\begin{enumerate}
\item Expand $\Lambda^{n-1}(\phi)(\epsilon_j)$, and prove that the coefficient of $\epsilon_i$ is $C_{ij}$. Interpret this as the wedge calculation from Problem~\ref{e:determinant}(c) with one row and one column omitted.
\item Group the terms of the permutation formula according to the row used in a fixed column $j$, and prove that $\det(A)=\sum_ia_{ij}C_{ij}$. Using transposition, derive the row expansion $\det(A)=\sum_ja_{ij}C_{ij}$ for a fixed row $i$. Check both formulas when $n=1$.
\item Write the $(i,j)$ entry of $A\adj(A)$ as $\sum_ka_{ik}C_{jk}$. For $i\ne j$, interpret it as a row expansion of the matrix obtained from $A$ by replacing row $j$ with row $i$. Deduce that $A\adj(A)=\det(A)I_n$, and use column replacement to prove the identity in the other order. Explain why the argument works over a commutative ring, even if that ring has zero divisors.
\item When $\det(A)\ne0$, deduce that $A^{-1}=\det(A)^{-1}\adj(A)$. Suppose $Ax=b$, and let $A_i(b)$ be the matrix obtained by replacing column $i$ of $A$ with $b$. Substitute $b=\sum_jx_j\operatorname{col}_j(A)$ and use multilinearity to prove \emph{Cramer's rule}:
\[
x_i=\frac{\det(A_i(b))}{\det(A)}.
\]
Why is the solution unique?
\item Now let $A$ be a square matrix over a commutative ring $R$. The adjugate identity still holds, but which condition on $\det(A)$ allows the formula for the inverse? Using the multiplicativity of the determinant, whose exterior proof also applies to finite free modules, show that $A$ is invertible if and only if $\det(A)$ is a unit of $R$. Give a $1\times1$ matrix over $\ZZ$ with nonzero determinant which is not invertible over $\ZZ$.
\end{enumerate}
\end{problems}

\hrulefill

Let $A\in\Mat_n(\KK)$ with $n\geq1$, and let $t$ be a formal variable. The matrix $tI_n-A$ has entries in the commutative ring $\KK[t]$, and its determinant
\[
p_A(t)\coloneqq\det(tI_n-A)
\]
is called the \emph{characteristic polynomial} of $A$. The permutation formula shows that it is monic of degree $n$. A change of basis replaces $tI_n-A$ by $P^{-1}(tI_n-A)P$, so multiplicativity shows that $p_A(t)$ depends only on the endomorphism, and not on its matrix.

Why should any polynomial in $A$ vanish at all? Since $\Mat_n(\KK)$ has dimension $n^2$, the matrices $I_n,A,\ldots,A^{n^2}$ are linearly dependent, so there is some nonzero polynomial $f$ with $f(A)=0$. However, this argument gives no preferred polynomial, and only a bound of $n^2$ on its degree. The following theorem provides a particular monic polynomial of degree $n$. In the language of Worksheet 1.2, with $t$ in place of $x$, it says that $p_A(t)$ acts as zero on the entire module $V_T$ defined by the operator $T(v)=Av$, even though $p_A(t)$ was constructed from a top exterior power.

\begin{theorem}[Cayley--Hamilton]
Every square matrix over $\KK$ satisfies its characteristic polynomial: if $p_A(t)=t^n+c_{n-1}t^{n-1}+\cdots+c_0$, then
\[
A^n+c_{n-1}A^{n-1}+\cdots+c_0I_n=0.
\]
\end{theorem}

The adjugate identity over $\KK[t]$ gives $(tI_n-A)\adj(tI_n-A)=p_A(t)I_n$. It is tempting to substitute $A$ for $t$ and declare the left side to be zero. There is a point to justify here: the coefficients of $\adj(tI_n-A)$ are matrices, and substituting a matrix for $t$ does not in general preserve products of polynomials with matrix coefficients. Instead we compare coefficients of $t$ and arrange the resulting identities so that they cancel. This gives a proof which uses only operations in the matrix ring $\Mat_n(\KK)$, as we check in the following exercise.

\hrulefill

\begin{problems}
\item\label{e:cayleyhamilton} \textbf{Cayley--Hamilton by Comparing Coefficients.} Throughout this exercise let $A\in\Mat_n(\KK)$ with $n\geq1$.
\begin{enumerate}
\item Use the permutation formula to prove that $p_A(t)$ is monic of degree $n$. Show that every entry of $\adj(tI_n-A)$ has degree at most $n-1$, and write
\[
\adj(tI_n-A)=B_{n-1}t^{n-1}+\cdots+B_1t+B_0,\qquad B_j\in\Mat_n(\KK).
\]
Explain why multiplying these expressions takes place in $\Mat_n(\KK[t])$, where $t$ commutes with every constant matrix.
\item Compare coefficients in $(tI_n-A)\adj(tI_n-A)=p_A(t)I_n$ to obtain
\[
\begin{gathered}
B_{n-1}=I_n,\qquad -AB_0=c_0I_n,\\
B_{j-1}-AB_j=c_jI_n\quad(1\leq j\leq n-1).
\end{gathered}
\]
Check these statements when $n=1$, where the last range is empty.
\item Multiply the identity indexed by $j$ on the left by $A^j$. Add the resulting identities, the identity $-AB_0=c_0I_n$, and $A^nB_{n-1}=A^n$. Write out which terms cancel, and conclude that $p_A(A)=0$. Did this argument require commuting $A$ past any $B_j$?
\item For a linear operator $\phi:V\to V$ on a finite-dimensional space, use its matrix in a basis to deduce that $p_{\phi}(\phi)=0$. Explain why this statement does not depend on the basis. If $V=0$, what are the characteristic polynomial of $\phi$ and its value at $\phi$?
\item Deduce that every power $A^m$ with $m\geq n$ is a linear combination of $I_n,A,\ldots,A^{n-1}$. If $c_0\ne0$, rearrange $p_A(A)=0$ to express $A^{-1}$ as a polynomial in $A$. Compare this with the adjugate formula for the inverse.
\end{enumerate}
\end{problems}

\hrulefill

We have now obtained determinant identities in two complementary ways. Exterior powers explain why the determinant is intrinsic and multiplicative, while coordinates describe the induced maps using minors and allow us to calculate the adjugate. The proof of the Cayley--Hamilton theorem then uses an identity over $\KK[t]$ to obtain an identity for an operator over $\KK$. We finish by putting these constructions together in a single example, including the cancellation of coefficients from Problem~\ref{e:cayleyhamilton}.

\hrulefill

\begin{problems}
\item\label{e:example} \textbf{An Explicit Determinant and Cayley--Hamilton Calculation.} Throughout this exercise work over $\QQ$, and let
\[
A=\begin{pmatrix}1&2\\3&4\end{pmatrix},\qquad\phi:\QQ^2\to\QQ^2,\quad v\longmapsto Av.
\]
\begin{enumerate}
\item Compute $\phi(e_1)\wedge\phi(e_2)$ as a multiple of $e_1\wedge e_2$, and read off $\det(A)$. Compute the cofactor matrix, $\adj(A)$, and both products in the adjugate identity.
\item Find $A^{-1}$, and use it to solve $Ax=(1,0)^{\mathsf T}$. Solve the same system using Cramer's rule, writing down each matrix $A_i(b)$.
\item Compute $p_A(t)$ and $\adj(tI_2-A)$. Check that $p_A(t)=t^2-5t-2$, and write the adjugate as $B_1t+B_0$. Verify each of the coefficient identities from Problem~\ref{e:cayleyhamilton}(b).
\item Compute $A^2$ and verify directly that $A^2-5A-2I_2=0$. Recover the same identity by the cancellation of Problem~\ref{e:cayleyhamilton}(c), and then use it to express $A^3$ and $A^{-1}$ as linear combinations of $I_2$ and $A$.
\item Every entry of $A$ and every coefficient of $p_A(t)$ is an integer. Reduce the Cayley--Hamilton identity modulo $2$, and check it for the resulting operator on $\FF_2^2$. Is the reduced matrix invertible? Which identities still hold, and which step in your calculation of the inverse is no longer available?
\end{enumerate}
\end{problems}

\hrulefill

\end{document}
